\chapter*{How to use this course}
\addcontentsline{toc}{chapter}{How to use this course}
\markboth{How to use this course}{How to use this course}

\section*{Who this is for}

This course is written for someone who knows \hi{nothing} about operating systems. Nothing
about assembly language. Nothing about how a computer starts up. If you can write a C program
that prints \cod{"Hello"}, you have enough to begin. If you have also met a pointer without
running away, even better.

What it does ask of you is a particular kind of patience. An operating system is a machine of
interlocking parts: almost none of them makes complete sense until you have seen the next one.
There will be moments where you read ``we will cover this in Chapter 24'' and have to accept it.
That is normal, and it is deliberate --- trying to explain everything at once is precisely why
most people give up.

\section*{What you are going to build}

The system studied here is called \hi{lytlnyblOS}. It is deliberately small: around 5,500
lines of C and assembly. But it is \emph{complete} in the way that matters. It boots from nothing
on a virtual machine, takes control of the processor, manages memory, runs several programs at
once, isolates them from each other, reads and writes a disk through its own filesystem, and
leaves you at a command line where you can type \cod{ls /bin} and get an answer.

None of that comes for free. There is no library underneath. There is no Linux underneath. When
the system prints a letter on screen, it is because someone wrote a byte to memory address
\hexa{B8000}. When it reads a key, it is because someone programmed the interrupt controller to
raise a signal. That ``someone'' is going to be you, and understanding exactly why each piece sits
where it does is the point of this document.

\begin{concept}{The journey in one sentence}
You will go from ``the computer powers on and knows how to do nothing'' to ``an unprivileged user
program asks the kernel to print text for it, and the kernel obliges'' --- understanding every
step in between.
\end{concept}

\section*{How this is organised}

The course has thirteen parts plus appendices. The order is \hi{not negotiable}: each part
builds on the previous one, exactly as the operating system itself does.

\begin{itemize}
  \item \hi{Parts I and II} are groundwork: what a processor is, what memory is, what an
        operating system actually is and which minimum pieces it needs. There is no project code
        yet, but nothing later stands up without this.
  \item \hi{Parts III to VII} build a minimal kernel: boot, enter protected mode, print to
        screen, handle interrupts, and talk to the timer and keyboard. By the end of Part VII you
        have a system that reacts to the outside world.
  \item \hi{Parts VIII and IX} are the heart of it: memory management and processes. This is
        where ``a program that runs'' becomes ``an operating system''.
  \item \hi{Parts X to XII} open the system up: user space, system calls, disk, filesystem
        and the shell.
  \item \hi{Part XIII} is real engineering: how it all compiles, how to debug it, and a
        complete worked case study of diagnosing a real fault in this project.
\end{itemize}

\section*{The boxes}

Five kinds of box appear throughout the text. It is worth knowing what each colour means.

\begin{concept}{Concept}
Blue. A central idea you will reuse many times. If you could only keep one thing from the page,
this is it.
\end{concept}

\begin{analogy}
Purple. A comparison with something from everyday life. Useful for building intuition, but
\emph{every analogy lies a little}: when intuition and code disagree, the code wins.
\end{analogy}

\begin{warn}{Watch out}
Amber. A classic mistake, a subtle trap, or something that costs hours of debugging in practice.
\end{warn}

\begin{danger}{This breaks the system}
Red. Serious faults: memory corruption, reboot loops, the kind of thing that does not produce an
error message but simply makes nothing work.
\end{danger}

\begin{tryit}
Green. Something you can run or look at right now to confirm that what you just read is true. Use
them: reading about operating systems without verifying anything is the best way to believe you
have understood something you have not.
\end{tryit}

\section*{What you need installed}

All development happens on Linux (on Windows, through WSL). You need five tools:

\begin{center}
\begin{tabularx}{\textwidth}{@{}l l X@{}}
\toprule
\textbf{Tool} & \textbf{Purpose} & \textbf{Why that one} \\
\midrule
\cod{nasm} & Assembler & Intel syntax, produces flat binaries with no header \\
\cod{i686-elf-gcc} & C compiler & A \emph{cross} compiler: generates code for a machine with no operating system \\
\cod{ld} & Linker & Places each chunk of code at the exact address we want \\
\cod{qemu-system-i386} & Emulator & Emulates a whole PC; if the system hangs, you close a window \\
\cod{gdb} & Debugger & Lets you stop the emulated processor and inspect its registers \\
\bottomrule
\end{tabularx}
\end{center}

The delicate one is the cross compiler. Your system's \cod{gcc} produces programs for \emph{your}
Linux: they expect a standard library, a dynamic loader, and a kernel to provide services. Our
operating system has none of that, because it \emph{is} the thing that would provide it. Hence a
\cod{gcc} configured for the \cod{i686-elf} target: ``32-bit x86 processor, ELF format, no
operating system behind it''. Part III explains this in detail.

\begin{tryit}[Check your environment]
\begin{lstlisting}[style=sh]
nasm -v
i686-elf-gcc --version
qemu-system-i386 --version
make --version
gdb --version
\end{lstlisting}
If all five answer, you can run \cod{make run} inside \filepath{lytlnybl/} and watch the system
boot.
\end{tryit}

\section*{An honest warning about the code}

This course explains a \emph{teaching} operating system. That means some decisions were made so
the code reads clearly, not so it is robust. At several points you will find red boxes pointing
at real limitations of the project: missing validation, memory that is never freed, edge cases
that hang the machine.

That is not a flaw in the course --- \hi{it is part of the course}. Learning where a real
operating system would have to be more careful is as valuable as learning how it behaves when
everything goes right. Part XIII collects all of those points and explains how each would be
fixed.

\sectionbreak
Let us start at the real beginning: what a computer is.
